\chapter{Chorionic Villus Sampling}

\section*{Overview}
Chorionic villus sampling (CVS) is a prenatal test that obtains a small sample of placental tissue, the chorionic villi, to detect chromosomal and genetic abnormalities in the fetus. It is typically performed between 10 and 13 weeks of gestation.

\section*{Alternative Names}
CVS, Chorionic villi biopsy, Placental biopsy

\section*{Principle}
A small sample of chorionic villi, which derive from the placenta and share the genetic makeup of the fetus, is obtained through the cervix or the abdominal wall under ultrasound guidance. The villi are cultured and analyzed for chromosomal karyotype, molecular genetic, or enzymatic studies. Because the villi and the fetus arise from the same zygote, results reflect fetal genetic material.

\section*{History}
Chorionic villus sampling was developed in the 1960s and 1970s and came into clinical use in the 1980s as an earlier alternative to amniocentesis.

\section*{Why It Is Done}
CVS is ordered for women at increased risk of fetal chromosomal or genetic abnormalities, such as those with advanced maternal age, abnormal prenatal screening results, or a family history of a genetic disorder. It provides earlier diagnosis than amniocentesis, allowing earlier decisions about the pregnancy.

\section*{Is This a Primary or Secondary Test or Procedure}
Primary invasive prenatal diagnostic test for genetic evaluation in the first trimester.

\section*{How It Is Performed}
Under continuous ultrasound guidance, a catheter or needle is introduced into the placenta either transvaginally or through the abdominal wall, and a small amount of chorionic villi is aspirated. The sample is examined under a microscope to confirm adequate tissue. The procedure takes several minutes and generally does not require anesthesia, though local anesthesia may be used for the abdominal route.

\section*{Preparation}
Informed consent is obtained and maternal blood type is reviewed so that anti-D immune globulin can be given to Rh-negative women. No fasting is required, and a full bladder may aid the transcervical approach.

\section*{Contraindications and Precautions}
The transcervical route is avoided with active vaginal bleeding, cervical infection, or cervical lesions. Relative contraindications include a posterior placenta, uterine fibroids, or an approach that is technically difficult. Vaginal spotting is common afterward and should be reported if heavy.

\section*{Risks}
Miscarriage occurs in approximately 0.5--1\% of cases, slightly higher than with amniocentesis. Cramping, vaginal spotting, leakage of fluid, and infection may occur, and limb reduction defects were reported when sampling was performed before 10 weeks.

\section*{Aftercare and Recovery}
The patient rests briefly and usually returns to normal activity the same day. Results are reported after culture or molecular analysis is complete, typically within one to two weeks.

\section*{Outcome and Follow-up}
Results report the fetal chromosomal pattern or the status of tested genes. Abnormal results are reviewed with a genetic counselor, while normal results reduce but do not eliminate the risk of birth defects.

\section*{Expected Results}
Normal karyotype (46,XX or 46,XY) with no evidence of the tested genetic disorders.

\section*{Follow-up}
Genetic counseling is provided for abnormal results, and a detailed fetal ultrasound evaluates for structural anomalies. Confirmatory amniocentesis may be offered when mosaicism or uncertain results are found.

\section*{When to Seek Medical Care}
Follow the procedure team's specific instructions. Seek urgent medical assessment for severe or worsening pain, uncontrolled bleeding, fever or signs of infection, trouble breathing, chest pain, fainting, new weakness or confusion, inability to urinate, or any rapidly worsening symptom. The appropriate warning signs depend on the procedure and the patient's medical condition.

\section*{References}
\begin{enumerate}
  \item MedlinePlus. Chorionic villus sampling. U.S. National Library of Medicine; 2025.
  \item Current Medical Diagnosis and Treatment. McGraw-Hill; 2025.
\end{enumerate}

\clearpage
